\clearpage\thispagestyle{empty}
\pdfbookmark[0]{Sources for the December 1804 Juno Report}{juno180412-sources}
\begin{center}\textbf{Sources for the December 1804 Juno Report}\end{center}
\noindent This AI translation follows Franz Xaver von Zach’s report as excerpted in \href{https://gdz.sub.uni-goettingen.de/id/PPN236007467}{Gauss’s \textit{Werke}, vol.~VI} (Göttingen, 1874), pp.~258--260. The corresponding passage of \textit{Ueber die Juno}, article LII of the \textit{Monatliche Correspondenz}, vol.~X, December 1804, begins on \href{https://www.digitale-sammlungen.de/en/view/bsb10538602?page=566}{p.~552} and reaches the Seeberg observations on \href{https://www.digitale-sammlungen.de/en/view/bsb10538602?page=569}{p.~555}.
\par\smallskip\noindent Von Zach quotes Carl Friedrich Gauss on the third set of elements. The tables include observations by von Zach at Seeberg and Nevil Maskelyne at Greenwich. The journal prints the comparison tables on \href{https://www.digitale-sammlungen.de/en/view/bsb10538602?page=567}{p.~553} and the ephemeris on \href{https://www.digitale-sammlungen.de/en/view/bsb10538602?page=568}{p.~554}. The collected edition describes the ephemeris in square brackets.
\par\smallskip\noindent Translation and typesetting produced by AI. Journal scans from the Bayerische Staatsbibliothek.
\clearpage
